\section{由赝 PDF 方法提取 K 介子胶子 PDF}\label{sec:results}

Ioffe 时间赝分布（pITD）~\cite{Radyushkin:2017cyf,Orginos:2017kos} 定义为
%
\begin{equation}
\mathcal{M}(\nu,z^2) = \langle 0 (P_z)|{\cal O}(z)|0 (P_z)\rangle,
\label{eq:ME_unpol}
\end{equation}
%
约化 pITD（RpITD）~\cite{Orginos:2017kos,Zhang:2018diq,Li:2018tpe} 的构造目的是消除 pITD 中的紫外发散：先取 pITD 与其在 $P_z=0$ 时相应的 $z$ 依赖矩阵元之比，再用 $z^2=0$ 处的矩阵元对该比值归一化：
 %
\begin{equation}
\mathscr{M}(\nu,z^2)=\frac{\mathcal{M}(zP_z,z^2)/\mathcal{M}(0\cdot P_z,0)}{\mathcal{M}(z\cdot 0,z^2)/\mathcal{M}(0\cdot 0,0)}.
\label{eq:RITD}
\end{equation}
%
在 RpITD 中，${\cal O}(z)$ 的重正化与运动学因子相消。
按此构造，本文采用的 RpITD 双比值在 $z=0$ 处归一化为 1。

图~\ref{fig:ensemble_momentum_comparison} 展示了 K 介子胶子 RpITD 对格距的依赖性。
图中 RpITD 表示为 Wilson 链长度 $z$ 的函数，系综 a12m310 与 a15m310（格距分别约为 0.12 与 0.15~fm）取近似相同的动量 $P_z$。
我们观察到，在每一个 Wilson 链长度处，a12m310 系综的 RpITD 数值都更高，且数值下降得更慢。
大多数数据点与另一系综的对应点相差在一个标准差之内，但多数 $z$ 处中心值略大。
随着 Wilson 链长度的增大，差异随之增大。


\begin{figure}[!tbp]
\centering
\includegraphics[width=0.45\textwidth]{images/ensemble_momentum_comparison.pdf}
\caption{\label{fig:ensemble_momentum_comparison}
a12m310 与 a15m310 系综 K 介子 RpITD 在助推动量 $P_z \approx 1.3$~GeV 下随 Wilson 线位移 $z$ 的变化。
}
\end{figure}


利用式~\ref{eq:3ptC} 与式~\ref{eq:RITD} 得到的矩阵元，我们计算了两个系综在不同 Wilson 链长度 $z$ 与介子动量 $P_z$ 处的 RpITD。
图~\ref{fig:RITD_fitbands} 将这些结果与文献~\cite{Fan:2021bcr} 此前在 a12m310 系综上计算的 $\pi$ 介子 RpITD 进行了比较。
可以看到，左图与中图所示两个系综的 K 介子 RpITD 相互吻合。
此外，最右图所示 a12m310 上 K 介子与 $\pi$ 介子的 RpITD 十分相近，这与文献~\cite{Roberts:2021nhw} 的结果一致——该文献计算得到 K 介子 PDF 与 $\pi$ 介子 PDF 之比约为 1。


\begin{figure*}[!tbp]
\centering
\includegraphics[width=0.32\textwidth]{images/Kaon_RITD_a15m310.pdf}\hfill
\includegraphics[width=0.32\textwidth]{images/Kaon_RITD_a12m310.pdf}\hfill
\includegraphics[width=0.32\textwidth]{images/Pion_RITD_a12m310.pdf}
\caption{\label{fig:RITD_fitbands}
a15m310（左）、a12m310（中）系综的 K 介子 RpITD，以及 a12m310 的 $\pi$ 介子 RpITD（右）。
各带是通过极小化式~\ref{eq:chi2-RITD-fit} 定义的 $\chi^2$ 对每个 $z$ 拟合胶子 PDF 所得的拟合带。}
\end{figure*}


随后，我们可以通过赝 PDF 匹配条件~\cite{Balitsky:2019krf} 提取胶子 PDF 分布；该条件把 RpITD $\mathscr{M}$ 与光锥胶子 PDF $g(x,\mu^2)$ 联系起来：
%
\begin{align}
\mathscr{M}(\nu,z^2)&=\int_0^1 dx \frac{xg(x,\mu^2)}{\langle x \rangle_g}R_{gg}(x\nu,z^2\mu^2),
\label{eq:matching-gg}
\end{align}
%
其中 $\mu$ 是 $\overline{\text{MS}}$ 方案下的重正化标度，
$\langle x \rangle_g=\int_0^1 dx \, x g(x,\mu^2)$ 是 K 介子中胶子的动量分数。
$R_{gg}$ 是胶子在胶子中的匹配核：
%
\begin{equation}
R_{gg}(y,z^2,\mu^2)
=\cos y -\frac{\alpha_s(\mu)}{2\pi}N_c
\left\{ \left[\ln\left(z^2\mu^2\frac{e^{2\gamma_E+1}}{4}\right)+2\right]R_B(y)+R_L(y)+R_C(y)\right\}, \nonumber
\label{eq:Rgg}
\end{equation}
%
其中 $\alpha_s$ 是标度 $\mu$ 处的强耦合常数，
$N_c=3$ 是颜色数，
$\gamma_E=0.5772$ 是欧拉—马歇罗尼（Euler--Mascheroni）常数。
对于项 $R_{gg}(y,z^2,\mu^2)$，按照此前关于赝 PDF 单圈演化的文献~\cite{Radyushkin:2018cvn}，$z$ 取为 $2e^{-\gamma_E-1/2}/\mu$ 以使对数项消失，从而压制含更高阶对数项的残余。
式~\ref{eq:matching-gg} 与诸项 $R_B(y)$、$R_L(y)$、$R_C(y)$ 的定义见文献~\cite{Balitsky:2019krf} 中的式 (7.21)--(7.23) 及式 (7.23) 之下的一段。

注意到，格点计算的 RpITD 还可通过夸克—胶子匹配核 $R_{gq}$ 与 K 介子的单态夸克 PDF $q_s$ 相联系：只需在式~\ref{eq:matching-gg} 中附加一项 $\frac{P_z}{P_0}\int_0^1 dx \frac{xq_S(x,\mu^2)}{\langle x \rangle_g}R_{gq}(x\nu,z^2\mu^2)$。
在此前的 $\pi$ 介子胶子 PDF 研究~\cite{Fan:2021bcr} 中发现夸克 PDF 贡献很小；
因此在本工作中我们将忽略 K 介子夸克 PDF。
通过式~\ref{eq:matching-gg} 的匹配条件拟合 RpITD 即可获得胶子 PDF $g(x,\mu^2)$；
HadStruc 合作组也使用过类似的流程~\cite{Joo:2019bzr,Joo:2020spy,HadStruc:2021wmh}。

为得到式~\ref{eq:matching-gg} 右端的胶子 PDF $g(x,\mu^2)$，我们采用如下受唯象学启发的形式：
%
\begin{align}
f_g(x,\mu) = \frac{xg(x, \mu)}{\langle x \rangle_g(\mu)} = \frac{x^A(1-x)^C}{B(A+1,C+1)},
\label{functional}
\end{align}
%
其适用于 $x\in[0,1]$，其余处为零。
贝塔函数 $B(A+1,C+1)=\int_0^1 dx\, x^A(1-x)^C$ 用于把面积归一化为 1。
JAM 在获取 $\pi$ 介子胶子 PDF 的全局拟合中也采用了同样的形式~\cite{Barry:2018ort,Cao:2021aci}。
%
我们通过极小化 $\chi^2$ 函数，将式~\ref{eq:RITD} 得到的格点 RpITD $\mathscr{M}^\text{lat} (\nu,z^2,a,M_\pi)$ 拟合到式~\ref{eq:matching-gg} 的参数化形式 $\mathscr{M}^\text{fit}(\nu,\mu,z^2,a,M_\pi)$：
%%
\begin{equation}
\label{eq:chi2-RITD-fit}
    \begin{split}
    & \chi^2(\mu,a,M_\pi)=\\
    & \sum _{\nu,z} \frac{(\mathscr{M}^{\rm fit}(\nu,\mu,z^2,a,M_\pi)-\mathscr{M}^{\rm lat}(\nu,z^2,a,M_\pi))^2}{\sigma^2_{\mathscr{M}}(\nu,z^2,a,M_\pi)}.
    \end{split}
\end{equation}
a15m310 与 a12m310 系综 K 介子 RpITD 以及 a12m310 的 $\pi$ 介子 RpITD 在各 $z^2$ 处的重建拟合带与格点计算点的比较，从左到右示于图~\ref{fig:RITD_fitbands}。
我们看到，a12m310 的两种情形中重建带几乎没有 $z^2$ 依赖性（以不同颜色标注），而 a15m310 情形的依赖稍强。
a15m310 K 介子数据的 RpITD 数据点涨落更大，因此拟合质量不如 a12m310 情形。
我们发现，a12m310 上对 $\pi$ 介子与 K 介子胶子 PDF 的拟合质量都非常稳定，$\chi^2/\text{dof}$ 约为 1，且无论 Wilson 线位移最大值 $z$ 如何选取，输出 $f_g(x,\mu)$ 都彼此一致。
然而对于 a15m310 系综，$\chi^2/\text{dof}$ 可高达 9(3)。
把参与拟合的 $z$ 减小到 $z=1$ 时，$\chi^2/\text{dof}$ 可降至 5(3)。
我们怀疑在此粗格距下高扭度效应被增强，致使拟合无法准确描述格点数据。
未来纳入 NNLO 匹配的工作或许有助于改善该系综上的拟合。



\begin{figure*}[!tbp]
  \centering
\includegraphics[width=0.45\textwidth]{images/KaonPDF_Lattice_DSE.pdf}\hfill
\includegraphics[width=0.45\textwidth]{images/xg-x0-1_pion_kaon_a12m310_comparison.pdf}
\caption{\label{fig:DSE_comparison}
%
左图：由格点数据拟合得到的 K 介子胶子 PDF $xg(x, \mu)/\langle x \rangle_g$ 随 $x$ 的变化，系综格距 $a\approx\{0.12,0.15\}$~fm（插图），$a\approx 0.12$~fm 处 $M_\pi\approx310$~MeV，并与 $\overline{\text{MS}}$ 方案下 $\mu=2$~GeV 的 DSE'20 K 介子胶子 PDF 比较。
%
右图：$\pi$ 介子与 K 介子胶子 PDF $xg(x, \mu)/\langle x \rangle_g$ 随 $x$ 变化的比较，格距 $a\approx0.12$~fm（外图）与 0.15~fm（插图），$M_\pi\approx310$~MeV，$\overline{\text{MS}}$ 方案、$\mu=2$~GeV。
}
\label{fig:xgx}
\end{figure*}


由格点数据拟合并表示为 $x$ 函数的 K 介子与 $\pi$ 介子胶子 PDF $xg(x, \mu)/\langle x \rangle_g$ 结果示于图~\ref{fig:xgx}。
左图插图中，我们给出了两个格距 $a\approx\{0.12,0.15\}$~fm、$M_\pi\approx 310$~MeV 在 $\overline{\text{MS}}$ 方案 $\mu=2$~GeV 处的 K 介子 $xg(x, \mu)/\langle x \rangle_g$。
注意，较粗格距我们采用 $z=1$，因为随着 $z$ 增大拟合质量显著变差；而对较细格点的 RpITD 数据，直到 $z=5$ 我们都能获得良好拟合。
采用这些数据输入选择，我们对 K 介子 $xg(x, \mu)/\langle x \rangle_g$ 的最佳测定在统计误差内相互一致。
我们还将较细格点的结果与由 Dyson--Schwinger 方程（DSE）方法~\cite{Cui:2020tdf} 得到的同一物理量进行了比较，后者在左图中以紫色带表示。
%
文献~\cite{Cui:2020tdf} 给出了近期关于 $\pi$ 介子胶子 PDF 以及 K 介子与 $\pi$ 介子胶子 PDF 之比随 $x$ 变化的 DSE 计算。
虽然其中未直接给出 K 介子胶子 PDF，但它可由 $\pi$ 介子胶子 PDF 与 K/$\pi$ 比值的乘积得到。
因此，我们可以将我们的格点计算与 DSE 计算的 K 介子胶子 PDF 进行直接比较。
另一方面，由于已发表的 K 介子与 $\pi$ 介子 PDF 参数缺乏相关性信息，图中所示 DSE'20 的误差很可能被高估。
%
我们在最小格距 $a\approx 0.12$~fm 处的 K 介子胶子 PDF 结果，在 $x>0.15$ 范围内与 DSE 结果~\cite{Cui:2020tdf} 在一个标准差误差内一致。
在图~\ref{fig:xgx} 的右侧，我们把较小格距得到的 K 介子结果与相同系综得到的 $\pi$ 介子结果进行比较；
可以看到，K 介子胶子 PDF 比 $\pi$ 介子的略小，这与其相应的价夸克 PDF 及 DSE~\cite{Cui:2020tdf} 结果类似。
%
我们预言 K 介子的 $\langle x^2 \rangle_g/\langle x \rangle_g$ 与 $\langle x^3 \rangle_g/\langle x \rangle_g$ 分别为 0.0779(94) 与 0.0187(42)，与 DSE~\cite{Cui:2020tdf} 的对应结果 0.075 与 0.015 符合良好。
对于 $\pi$ 介子的 $\langle x^2 \rangle_g/\langle x \rangle_g$，我们的 310~MeV 结果给出
0.092(15) 与 0.0250(75)，
而 DSE~\cite{Cui:2020tdf}、JAM~\cite{Barry:2018ort,Cao:2021aci} 与 xFitter~\cite{Novikov:2020snp} 的结果分别为
0.076、0.103、0.158
与 0.015、0.024、0.048。
%
纳入更细格距与更轻 $\pi$ 介子质量的未来研究，将对完善这一计算十分重要，并为这个知之甚少的介子物理量提供更好的预言。
